\documentclass[a4paper,10pt]{article}
\hyphenpenalty=10000
\exhyphenpenalty=10000

% === Packages ===
\usepackage[left=0.7in, right=0.7in, top=0.3in, bottom=0.3in]{geometry}
\usepackage{xcolor}
\usepackage{titlesec}
\usepackage{enumitem}
\usepackage{fontspec}
\newfontfamily\datefont{Latin Modern Roman}
\usepackage{multicol}
\usepackage{fancyhdr}
\usepackage{tabularx}
\usepackage{hyperref}
\usepackage{fontawesome5}

% === Font ===
\setmainfont{Inconsolata}

% === Spacing ===
\linespread{0.90}
\setlength{\parskip}{0pt}
\setlength{\parindent}{0pt}
\setlength{\columnsep}{16pt}
\setlength{\multicolsep}{2pt}

\setlist[itemize]{
  leftmargin=*,
  itemsep=1.5pt,
  topsep=2pt,
  parsep=0pt,
  partopsep=0pt
}

% === Section styling ===
\definecolor{sectioncolor}{RGB}{40,70,140}
\titleformat{\section}
  {\large\bfseries\color{sectioncolor}}
  {}
  {0em}
  {}
  [\titlerule]

\titlespacing*{\section}{0pt}{4pt}{2pt}

% === Hyperref ===
\hypersetup{
  colorlinks=true,
  urlcolor=sectioncolor,
  linkcolor=sectioncolor
}

% === Header/Footer ===
\pagestyle{fancy}
\fancyhf{}
\renewcommand{\headrulewidth}{0pt}
\renewcommand{\footrulewidth}{0pt}

\begin{document}
\sloppy

% === HEADER ===
\begin{center}
  {\Huge Ipsit Mohanty} \\[2.5mm]
  {\Large Data \& ML Engineer} \\[1mm]
  {\small \textit{Python • ETL Pipelines • Analytical Warehousing • Machine Learning • FastAPI • GenAI (RAG)}} \\[2mm]
  \faEnvelope~\href{mailto:ipsit13@gmail.com}{ipsit13@gmail.com} \quad
  \faPhone~8433904389 \quad
  \faLinkedin~\href{%%TRACKER_LI%%}{LinkedIn} \quad
  \faGithub~\href{%%TRACKER_GH%%}{GitHub}
\end{center}
\noindent\rule{\textwidth}{0.2pt}

% === SUMMARY ===
\section*{Summary}
Data and ML engineer with 6+ years building ETL pipelines, analytical warehouses, and machine
learning systems on large-scale government and operational datasets spanning 70,000+ entities. Ships
end-to-end: ingestion, validation, and schema design through models, FastAPI services, dashboards,
and automated reporting. Recent hands-on work spans deep learning (CNN and vision-transformer
architectures in Keras and PyTorch) and GenAI (retrieval-augmented assistants with vector search
and guarded LLM fallback). Focused on reproducible data engineering, applied machine learning, and
operational decision-support systems.

% === EXPERIENCE ===
\section*{Experience}

% --- DIC ---
{\large\textbf{Digital India Corporation (NeGD), MeitY}} \\
\textit{State Coordinator - Data Systems \& Analytics | Poshan Tracker, Odisha} \hfill {\datefont\small Aug 2025 - Jun 2026}

\begin{itemize}
\item Engineered Python ETL pipelines for ingestion, cleaning, reconciliation, and analytics cube generation across 72,000+ Anganwadi centres in 30 districts, covering registration, supplementary nutrition, growth monitoring and health indicators for children under six and pregnant and lactating mothers.
\item Designed analytical warehouses (DuckDB/SQLite OLAP stores) with schema harmonisation across changing monthly export formats, anomaly and risk scoring, and month-on-month drift detection, replacing manual district-level reconciliation with automated validation layers.
\item Built the operational monitoring dashboard used in the monthly state review cycle, surfacing coverage, measurement-efficiency and nutrition indicators at state, district and centre level alongside anomaly and risk surveillance. A synthetic-data reconstruction is published as \href{%%TRACKER_REPO_AWC%%}{AWC Operations Dashboard}.
\item Built a RAG FAQ assistant for operational query resolution: deterministic FAQ matching, vector search, acronym expansion, and guarded LLM fallback with iteratively tuned prompt templates.
\item Designed FastAPI services exposing regional insights and cube access as structured analytical responses; automated PPT/PDF reporting and BI-ready extracts for state and district programme officers.
\item Acted as the technical interface between central and state teams, supporting implementation, troubleshooting, and adoption of data systems across distributed field operations.
\end{itemize}

\vspace{1pt}

% --- Piramal ---
\noindent\rule{\textwidth}{0.15pt}
{\large\textbf{Office of the Hon. Minister, Govt. of Odisha}} \\
\textit{Program Manager | Policy Analytics \& Systems Integration} \hfill {\datefont\small Dec 2024 - Aug 2025} \\
{\small \textit{Supported by Piramal Foundation}}

\begin{itemize}
\item Built dropout-risk analysis on UDISE+ national school enrolment data, combining clustering of school-level indicators with time-series trend analysis to support education-sector planning inputs under the National Education Policy.
\item Ran root-cause analytics on learning outcomes, dropout clusters, absenteeism, and teacher-deficit mapping; findings fed strategy briefs presented to department leadership.
\item Designed and built a public grievance workflow end-to-end (Flask/FastAPI): automated classification, resolution tracking, and executive reporting.
\item Implemented financial analytics for variance analysis, forecasting, and budget-utilisation monitoring on a major state education programme.
\end{itemize}

\vspace{1pt}

% --- STSC ---
\noindent\rule{\textwidth}{0.15pt}
{\large\textbf{ST \& SC Development Dept., Govt. of Odisha}} \\
\textit{Consultant | Analytics \& Governance Monitoring} \hfill {\datefont\small Apr 2023 - Nov 2024}

\begin{itemize}
\item Built the monitoring dashboard tracking 6,000+ cultural-infrastructure projects across 30 districts: fund disbursement, milestone adherence, and physical-progress reporting for district administrations.
\item Designed fund-flow analytics linking physical progress to allocation patterns, surfacing where projects were stalling and which bottlenecks were driving it, so remedial action could be targeted.
\item Applied clustering to prioritise high-need intervention zones, feeding district allocation reviews over a programme that expanded from 8 to 14 districts during the engagement.
\item Developed forecasting models (ARIMA, Prophet, regression trend analysis) for resource planning inputs in programme review cycles.
\end{itemize}

\vspace{1pt}

% --- PHDMA ---
\noindent\rule{\textwidth}{0.15pt}
{\large\textbf{Planning \& Convergence Dept., Govt. of Odisha}} \\
\textit{Young Professional | Governance \& Data Analytics} \hfill {\datefont\small Dec 2020 - Mar 2023}

\begin{itemize}
\item Built impact-evaluation models (gradient boosting) assessing programme delivery across six departments (agriculture, education, fisheries, revenue, home, and women's economic empowerment), reconciling field reporting from separately-executed programmes into comparable measures for fund-allocation reviews.
\item Automated variance tracking in financial disbursements, replacing manual reconciliation, and designed an anomaly-detection approach for early identification of irregularities in fund distribution.
\item Led sentiment analysis of citizen feedback at scale, supporting responsiveness under state governance and service-delivery initiatives.
\end{itemize}

% === SKILLS ===
\section*{Skills \& Methodologies}

\begin{multicols}{2}
\raggedright

\textbf{Data Engineering \& Analytics}
\begin{itemize}
\item ETL/ELT pipelines, schema design, validation, analytical warehousing
\item MySQL, PostgreSQL, DuckDB, SQLite, Apache Spark, Pandas, NumPy
\item Streamlit, Power BI, automated reporting
\end{itemize}

\textbf{Machine Learning}
\begin{itemize}
\item Supervised and unsupervised ML, feature engineering
\item Deep learning: CNNs, vision transformers, transfer learning, ONNX
\item Forecasting (ARIMA, Prophet, LSTM), anomaly detection; Scikit-learn, TensorFlow/Keras, PyTorch
\end{itemize}

\columnbreak

\textbf{Cloud, Backend \& MLOps}
\begin{itemize}
\item Google Cloud: BigQuery, BigQuery ML, Dataplex, GKE
\item FastAPI, Flask, REST services; Docker, Git, CI/CD (GitHub Actions), pytest, Linux
\item MLflow, DVC, experiment tracking
\end{itemize}

\textbf{GenAI}
\begin{itemize}
\item RAG pipelines, vector search, embeddings
\item Chroma, Sentence Transformers, Ollama
\item Prompt engineering, guarded LLM fallback design
\end{itemize}

\end{multicols}

% === PROJECTS ===
\section*{Selected Projects}

\noindent{\small All repositories at \textbf{\href{%%TRACKER_GH%%}{github.com/IpsitMohanty}}: tested, documented, CI where noted; live demo linked.}\\[1mm]

\begin{itemize}[itemsep=0pt]
\item \textbf{\href{%%TRACKER_REPO_CNN%%}{Satellite Land Classification (CNN \& CNN-ViT)}:} Deep-learning pipeline built twice, in Keras and PyTorch, spanning data loading, CNN classifiers, and CNN-ViT hybrids via transfer learning, reproduced locally on CPU. A four-experiment controlled ablation isolated a Keras/PyTorch BatchNorm momentum convention mismatch (0.99 vs 0.1, same parameter name) as the root cause of a 70\% to 98.7\% accuracy gap between otherwise identical models; the hybrid was shown not to earn its complexity (0.58pp lower accuracy at 2x parameters, 2.6x training time). pytest suite (55 tests), CI, ONNX-backed \href{%%TRACKER_DEMO_CNN%%}{live demo}.
\item \textbf{\href{%%TRACKER_REPO_RAG%%}{RAG Ingestion Evaluation}:} LangChain retrieval pipeline over a deliberately heterogeneous corpus (app FAQ pairs and a nutrition-scheme policy PDF), methodology fixed before any run. A 51-query evaluation across 24 ingestion/chunking/retriever configurations found that similarity distances cannot serve as a confidence gate: unanswerable-query scores overlap answerable ones in every configuration, so no threshold separates good retrieval from bad. A follow-up arm added an LLM abstention judge on LangGraph, testing whether a content-based judge can gate what the score cannot, prediction pre-registered. Extended into a single-agent corrective RAG loop (LangGraph): self-graded retrieval, corrective re-retrieval, and faithfulness-checked generation with abstention; a 545-call evaluation vs.\ the non-agentic baseline rescued failures it structurally could not (0 to 41.7\% recall on hardest queries), zero fabrication, negatives reported openly. Local embeddings, no API key. pytest suite (168 tests), CI, \href{%%TRACKER_DEMO_RAG%%}{live demo}.
\item \textbf{\href{%%TRACKER_REPO_REC%%}{Course Recommender System}:} Comparative study of nine recommendation approaches (content-based, collaborative filtering, neural embeddings) on a 34K-user dataset, unified under a common \texttt{Recommender} interface and evaluated on a reproducible, seeded held-out ranking protocol (Precision@k / Recall@k). pytest suite (63 tests); deployed as an interactive Streamlit app. Evaluation surfaces an honest result, with only a clustering model beating a popularity baseline: \href{%%TRACKER_DEMO_REC%%}{live demo}.
\item \textbf{\href{%%TRACKER_REPO_POSHAN%%}{Poshan Intelligence Pipeline}:} End-to-end data platform with 12 ETL modules for distinct government report types, a district-level analytics cube, and RandomForest models predicting low-birth-weight and stunting rates, served via FastAPI. Dockerised, with GitHub Actions CI running pytest (29 tests), flake8, and mypy.
\item \textbf{\href{%%TRACKER_REPO_AWC%%}{AWC Operations Dashboard}:} Data mart with staging-to-fact-table design, anomaly/risk/alert scoring layers, schema-drift detection across monthly exports, and a dual-backend Streamlit dashboard (PostgreSQL/SQLite) built on population-weighted rate calculations. Runs on a synthetic generator with injected ground-truth anomalies, with detection recall reported openly. Deployed with managed PostgreSQL (Neon): \href{%%TRACKER_DEMO_AWC%%}{live demo}.
\end{itemize}

% === CERTIFICATIONS ===
\section*{Certifications \& Affiliations}

\begin{itemize}[itemsep=1pt]
\item IBM Professional Certificates (2024-2025): AI Engineering, Generative AI Engineering with LLMs, Machine Learning, Data Science, Applied DevOps Engineering
\item Google Cloud skill badges (2025-2026): BigQuery ML, BigQuery warehousing and streaming analytics, Dataplex, ML APIs, GKE; Member, Google Cloud Innovators Program
\end{itemize}

% === EDUCATION ===
\section*{Education}

\renewcommand{\arraystretch}{0.9}
\begin{tabularx}{\textwidth}{X r}

\textbf{International Institute of Information Technology, Bangalore}
& {\datefont\small Mar 2026-Present} \\
Executive Diploma in Machine Learning \& AI \\

\textbf{Indian Institute for Human Settlements, Bangalore}
& {\datefont\small 2019-2020} \\
Urban Fellows Programme \\

\textbf{Tata Institute of Social Sciences, Mumbai}
& {\datefont\small 2017-2019} \\
Master's in Disaster Management \\

\textbf{Indira Gandhi National Open University, New Delhi}
& {\datefont\small 2013-2016} \\
BA (Major) Psychology \\

\end{tabularx}

\end{document}
